\section{Related Work}
\label{sec:related}

\paragraph{Real charging datasets.}
ACN-Data \citep{lee2019acn} remains the reference workplace-charging corpus:
per-session records from Caltech, JPL, and an office site, including
driver-entered energy requests. ElaadNL publishes anonymized public-charging
transactions from the Netherlands \citep{elaadnl2020}, and the DOE-funded
EV~WATTS program released a large multi-venue evaluation dataset
\citep{evwatts2023}. The INL EV Project \citep{inl2015evproject} documents an
earlier (2011--2013) residential fleet. These corpora anchor our calibration
(\S\ref{sec:calibration}); their limitations---few sites, fixed
infrastructure, no SoC channel, restricted or aggregated building-side
data---motivate generation rather than reuse.

\paragraph{Building-energy data and simulation.}
On the building side, large public resources exist: NREL's ComStock
\citep{parker2023comstock} provides calibrated end-use load profiles for the
U.S. commercial stock, and the Building Data Genome~2 project
\citep{miller2020bdg2} publishes meter data from over a thousand buildings.
Physics-based simulation via EnergyPlus \citep{crawley2001energyplus} with the
DOE prototype building models \citep{pnnl_prototypes} is the standard for
controlled building-load studies. We use simulation as the generative
mechanism and the public corpora as \emph{ground truth for validation}, which
keeps the two roles---generation and evaluation---separate by construction.

\paragraph{Synthetic EV-charging data.}
Session-level generative models appear in charging-infrastructure simulators:
ACN-Sim \citep{lee2021acnsim} includes generative session models fitted to
ACN-Data, and ElaadNL distributes simulation profiles derived from its
transactions. These target the charging channel in isolation. General-purpose
deep generative models for tabular and time-series data---GAN-based
\citep{yoon2019timegan,esteban2017rgan} and diffusion-based more
recently---can imitate a single site's joint distribution, but they offer no
scenario control (one cannot move the learned site to another climate or
double its fleet), no mechanism for physical consistency across coupled
channels, and no per-value auditability. \generator{} takes the opposite
trade: parametric distribution families joined by an explicit copula, fitted
per behavioral region, embedded in a physical simulation---less expressive
than a deep density model at any single site, but controllable, auditable,
physically coherent across channels, and calibrated across multiple corpora.
We quantify what this trade costs (and buys) in
\S\ref{sec:representativeness}--\S\ref{sec:utility}.

\paragraph{Evaluating synthetic data.}
We follow the two-sided evaluation that has become standard: distributional
fidelity against held-out real data, and downstream utility via
train-on-synthetic/test-on-real (TSTR) protocols introduced for medical
time series \citep{esteban2017rgan}. For building-load fidelity we adopt the
calibration metrics of ASHRAE Guideline~14 \citep{ashrae2014g14}
($\cvrmse{}$, $\nmbe{}$), which are the accepted standard for comparing
simulated to metered building energy. Documentation follows Datasheets for
Datasets \citep{gebru2021datasheets} and the Croissant machine-readable
metadata format \citep{akhtar2024croissant}.
